\chapter{Rewriting and Normalization}
\label{ch:rewriting}

Rewriting is the second great engine of automated reasoning, and it differs from resolution in a way that matters for this book. Resolution derives new facts and keeps all of them until a contradiction appears. Rewriting \emph{replaces}: a term is repeatedly simplified, the earlier forms are discarded, and what remains is a canonical representative. A rewriting system is thus a machine that turns a family of equal expressions into a single one, and the structure it forgets in doing so can be stated exactly. This chapter proves the basic theory from the ground up, namely confluence, Newman's lemma and termination by interpretation, cites the critical pair lemma and the correctness of completion, and ends with the observation that a normal form is the extreme case of a compression that is sufficient for one task and insufficient for others.

\section{Abstract reduction systems}
\label{sec:ars}

An \emph{abstract reduction system} is a pair $(A,\to)$ with ${\to}\subseteq A\times A$. Write $\to^{+}$ and $\to^{*}$ for the transitive and reflexive-transitive closures, $\leftarrow$ for the converse, and $\leftrightarrow^{*}$ for the equivalence generated by $\to$. Elements $a,b$ are \emph{joinable}, written $a\downarrow b$, if $a\to^{*}c\leftarrow^{*}b$ for some $c$. An element is a \emph{normal form} if no step leaves it, and $b$ is a normal form \emph{of} $a$ if $a\to^{*}b$ and $b$ is a normal form.

\begin{definition}[Termination and confluence]
\label{def:conf}
The system $(A,\to)$ is
\begin{itemize}[nosep]
\item \emph{terminating} if there is no infinite sequence $a_0\to a_1\to a_2\to\cdots$;
\item \emph{confluent} if $b\leftarrow^{*}a\to^{*}c$ implies $b\downarrow c$;
\item \emph{locally confluent} if $b\leftarrow a\to c$ implies $b\downarrow c$;
\item \emph{Church--Rosser} if $a\leftrightarrow^{*}b$ implies $a\downarrow b$;
\item \emph{convergent} if it is terminating and confluent.
\end{itemize}
\end{definition}

\begin{lemma}[Confluence is the Church--Rosser property]
\label{lem:cr}
$(A,\to)$ is confluent if and only if it is Church--Rosser.
\end{lemma}

\begin{proof}
($\Leftarrow$) If $b\leftarrow^{*}a\to^{*}c$ then $b\leftrightarrow^{*}c$, so $b\downarrow c$.

($\Rightarrow$) Suppose $a\leftrightarrow^{*}b$, by induction on the length of a zigzag from $a$ to $b$. Length zero is trivial. Otherwise $a\leftrightarrow^{*}b'$ by a shorter zigzag, and $b'\to b$ or $b'\leftarrow b$. By the induction hypothesis $a\to^{*}c\leftarrow^{*}b'$. If $b\to b'$ then $b\to b'\to^{*}c$, so $a\to^{*}c\leftarrow^{*}b$. If $b'\to b$ then $c\leftarrow^{*}b'\to b$, and confluence gives $c\to^{*}d\leftarrow^{*}b$, so $a\to^{*}d\leftarrow^{*}b$.
\end{proof}

\begin{theorem}[Newman's lemma, 1942]
\label{thm:newman}
A terminating and locally confluent system is confluent.
\end{theorem}

\begin{proof}
Since the system is terminating, $\to^{+}$ is well-founded and we may argue by well-founded induction. Let $P(a)$ say that $b\leftarrow^{*}a\to^{*}c$ implies $b\downarrow c$. Assume $P(a')$ for every $a'$ with $a\to^{+}a'$, and let $b\leftarrow^{*}a\to^{*}c$. If $a=b$ then $b=a\to^{*}c$ so $b\downarrow c$, and symmetrically if $a=c$. Otherwise $a\to a_1\to^{*}b$ and $a\to a_2\to^{*}c$, and by local confluence $a_1\to^{*}d\leftarrow^{*}a_2$.

The induction hypothesis holds at $a_1$ and at $a_2$, since $a\to a_1$ and $a\to a_2$. Applying it at $a_1$ to $b\leftarrow^{*}a_1\to^{*}d$ gives $b\to^{*}e\leftarrow^{*}d$. Applying it at $a_2$ to $d\leftarrow^{*}a_2\to^{*}c$ gives $d\to^{*}f\leftarrow^{*}c$. Now $a\to a_1\to^{*}d$ gives $a\to^{+}d$, so the induction hypothesis holds at $d$, and applying it to $e\leftarrow^{*}d\to^{*}f$ gives $e\to^{*}g\leftarrow^{*}f$. Then $b\to^{*}e\to^{*}g$ and $c\to^{*}f\to^{*}g$, so $b\downarrow c$.
\end{proof}

\begin{corollary}[Unique normal forms]
\label{cor:nf}
Let $(A,\to)$ be convergent. Every $a\in A$ has exactly one normal form $\NF(a)$, and $a\leftrightarrow^{*}b$ if and only if $\NF(a)=\NF(b)$. If in addition the successors of any element are computable, then $\leftrightarrow^{*}$ is decidable.
\end{corollary}

\begin{proof}
A normal form exists by termination, since otherwise one could step forever. If $b,c$ are normal forms of $a$ then $b\leftarrow^{*}a\to^{*}c$, so $b\to^{*}d\leftarrow^{*}c$ by confluence; as $b,c$ admit no steps, $b=d=c$.

If $a\leftrightarrow^{*}b$ then Lemma~\ref{lem:cr} gives $a\to^{*}e\leftarrow^{*}b$, hence $\NF(a)=\NF(e)=\NF(b)$ by uniqueness. Conversely, $a\to^{*}\NF(a)=\NF(b)\leftarrow^{*}b$ gives $a\leftrightarrow^{*}b$. For the last claim, compute $\NF(a)$ and $\NF(b)$ by following any sequence of steps, which halts by termination, and compare.
\end{proof}

Corollary~\ref{cor:nf} is the abstract content of every canonizer. Note what it does \emph{not} say: all maximal reduction sequences from $a$ end at the same element, but nothing is claimed about their lengths or their contents.

\begin{proposition}[Strategy independence of the output]
\label{prop:strategy-independence}
In a convergent system, every maximal reduction sequence from $a$ is finite and ends at $\NF(a)$. In particular every strategy for choosing the next step, adaptive or not, has the same output, and strategies differ only in the traces they produce.
\end{proposition}

\begin{proof}
A maximal sequence is finite by termination and ends in a normal form of $a$, which is $\NF(a)$ by Corollary~\ref{cor:nf}.
\end{proof}

\begin{example}[Same output, different traces]
\label{ex:strategies}
Take rules $a\to b$ and $h(x)\to d(x,x)$ over a signature with constants $a,b$, a unary $h$ and a binary $d$. The rules have no overlap, and the system is terminating. Reducing $h(a)$ innermost gives $h(a)\to h(b)\to d(b,b)$ in two steps; reducing outermost gives $h(a)\to d(a,a)\to d(b,a)\to d(b,b)$ in three. The outputs coincide, in agreement with the proposition, while both the length and the rules used differ.
\end{example}

Proposition~\ref{prop:strategy-independence} places rewriting at the opposite pole from the search of Chapter~\ref{ch:search}. There, strategy affects which derivation is found and whether one is found at all. Here, in a convergent system, strategy affects only cost. This is the formal reason that rewriting-based canonization is a safe component of a larger prover: it is a function, not a search.

\section{Term rewriting and termination by interpretation}
\label{sec:trs}

Fix a signature as in Section~\ref{sec:unification} of Chapter~\ref{ch:resolution}. A \emph{rewrite rule} is a pair $l\to r$ of terms with $l$ not a variable and $\mathrm{Var}(r)\subseteq\mathrm{Var}(l)$. A \emph{term rewriting system} (TRS) $\mathcal{R}$ is a finite set of rules. The rewrite relation is
\[
s\to_{\mathcal{R}}t\quad\Longleftrightarrow\quad s|_p=\sigma(l)\ \text{and}\ t=s[\sigma(r)]_p\ \text{for some rule } l\to r,\ \text{position } p,\ \text{substitution }\sigma,
\]
where $s|_p$ is the subterm at position $p$ and $s[u]_p$ is $s$ with that subterm replaced by $u$.

The most direct method for proving termination assigns each symbol a function into a well-founded set.

\begin{proposition}[Termination by monotone interpretation]
\label{prop:interp}
Let each $n$-ary symbol $f$ be interpreted by a function $[f]:\N^n\to\N$ that is strictly increasing in each argument. For an assignment $\alpha:\mathcal{V}\to\N$ define $[t]_\alpha$ by $[x]_\alpha=\alpha(x)$ and $[f(t_1,\dots,t_n)]_\alpha=[f]([t_1]_\alpha,\dots,[t_n]_\alpha)$. If $[l]_\alpha>[r]_\alpha$ for every rule $l\to r$ and every assignment $\alpha$, then $\to_{\mathcal{R}}$ is terminating.
\end{proposition}

\begin{proof}
We show that $s\to_{\mathcal{R}}t$ implies $[s]_\alpha>[t]_\alpha$ for every $\alpha$. For a substitution $\sigma$ and assignment $\alpha$ put $\alpha_\sigma(x)=[\sigma(x)]_\alpha$; by induction on terms, $[\sigma(u)]_\alpha=[u]_{\alpha_\sigma}$. Hence for a step at the root, $s=\sigma(l)$ and $t=\sigma(r)$, we get $[s]_\alpha=[l]_{\alpha_\sigma}>[r]_{\alpha_\sigma}=[t]_\alpha$. For a step at a deeper position, induct on the depth of $p$: if $s=f(s_1,\dots,s_i,\dots,s_n)$ and the step takes place inside $s_i$, giving $s_i'$ with $[s_i]_\alpha>[s_i']_\alpha$, then strict monotonicity of $[f]$ in its $i$-th argument gives $[s]_\alpha>[t]_\alpha$. An infinite reduction $s_0\to s_1\to\cdots$ would give, for any fixed $\alpha$, an infinite strictly decreasing sequence $[s_0]_\alpha>[s_1]_\alpha>\cdots$ in $\N$, which is impossible.
\end{proof}

\begin{example}[Addition]
\label{ex:addition}
Let $\mathcal{R}_{+}$ consist of $\mathit{add}(0,y)\to y$ and $\mathit{add}(s(x),y)\to s(\mathit{add}(x,y))$. Put $[0]=1$, $[s](x)=x+1$, $[\mathit{add}](x,y)=2x+y$. All three are strictly increasing in each argument. For the first rule, $[\mathit{add}(0,y)]_\alpha=2+\alpha(y)>\alpha(y)$. For the second, $[\mathit{add}(s(x),y)]_\alpha=2(\alpha(x)+1)+\alpha(y)=2\alpha(x)+\alpha(y)+2$, which exceeds $[s(\mathit{add}(x,y))]_\alpha=2\alpha(x)+\alpha(y)+1$. By Proposition~\ref{prop:interp}, $\mathcal{R}_{+}$ terminates. The two left-hand sides have no overlap, so there are no critical pairs in the sense of the next section, and $\mathcal{R}_{+}$ is convergent.
\end{example}

\section{Critical pairs and local confluence}
\label{sec:critical}

Because termination is available by Proposition~\ref{prop:interp} and Theorem~\ref{thm:newman} reduces confluence to a local condition, what remains is to check local confluence of a TRS. The check is finite, and it is carried out by unification.

\begin{definition}[Critical pair]
\label{def:critical-pair}
Let $l_1\to r_1$ and $l_2\to r_2$ be rules renamed to share no variables, and let $p$ be a position of $l_1$ such that $l_1|_p$ is not a variable and unifies with $l_2$ with most general unifier $\sigma$. If $p$ is the root then the two rules must not be variants of each other. The pair
\[
\langle\ \sigma(r_1),\ \sigma(l_1[r_2]_p)\ \rangle
\]
is a \emph{critical pair}. It is \emph{joinable} if its two components are joinable under $\to_{\mathcal{R}}$.
\end{definition}

\begin{lemma}[Critical pair lemma]
\label{lem:cpl}
\cited{Knuth--Bendix 1970, Huet 1980} A TRS is locally confluent if and only if all its critical pairs are joinable.
\end{lemma}

Because a finite TRS has finitely many critical pairs, each computed by the unification procedure of Theorem~\ref{thm:unification}, the lemma yields a decision procedure for confluence of terminating systems.

\begin{proposition}[Confluence of terminating TRS is decidable]
\label{prop:conf-decidable}
There is an algorithm that, given a finite terminating TRS, decides whether it is confluent.
\end{proposition}

\begin{proof}
Compute the finite set of critical pairs. For each pair, enumerate the reducts of each component. In a terminating, finitely branching system the set of reducts of a term is finite: otherwise by K\"onig's lemma the finitely branching tree of reduction sequences would contain an infinite branch, contradicting termination. So joinability is decidable by checking whether the two finite reduct sets intersect. By Lemma~\ref{lem:cpl}, the system is locally confluent exactly when all pairs are joinable, and by Theorem~\ref{thm:newman} local confluence is confluence for terminating systems.
\end{proof}

\begin{example}[Groups]
\label{ex:groups}
Over the signature $\{\cdot,{}^{-1},e\}$ consider the ten rules
\begin{center}
\begin{tabular}{lll}
$e\cdot x\to x$ & $x^{-1}\cdot x\to e$ & $(x\cdot y)\cdot z\to x\cdot(y\cdot z)$\\
$x^{-1}\cdot(x\cdot y)\to y$ & $x\cdot e\to x$ & $e^{-1}\to e$\\
$(x^{-1})^{-1}\to x$ & $x\cdot x^{-1}\to e$ & $x\cdot(x^{-1}\cdot y)\to y$\\
$(x\cdot y)^{-1}\to y^{-1}\cdot x^{-1}$
\end{tabular}
\end{center}
This is the classical completion of the three group axioms, due to Knuth and Bendix. Termination follows from the lexicographic path order with precedence $\mathord{}^{-1}>\cdot>e$, a standard result that we cite rather than reprove. Local confluence was checked in preparation of this book by computing every critical pair over all ordered pairs of rules and all non-variable positions, 55 in total, and normalizing each side: all were joinable. This is a computation by a short independent program and has not been mechanized in a proof assistant. Granting termination, Proposition~\ref{prop:conf-decidable} then makes the system convergent, so by Corollary~\ref{cor:nf} two group terms are equal in every group if and only if they have the same normal form (using the equational completeness theorem recalled below).
\end{example}

\section{Equational theories and completion}
\label{sec:completion}

Let $E$ be a set of equations $s\approx t$ between terms. The relation $\leftrightarrow_E$ is the symmetric closure of single rewrite steps using equations of $E$ in both directions, at any position and under any substitution, and $\leftrightarrow_E^{*}$ is its reflexive-transitive closure. Write $E\models s\approx t$ if every algebra satisfying all equations of $E$ satisfies $s\approx t$ under every assignment.

\begin{theorem}[Birkhoff 1935]
\label{thm:birkhoff}
\cited{Birkhoff} $E\models s\approx t$ if and only if $s\leftrightarrow_E^{*}t$.
\end{theorem}

Birkhoff's theorem reduces equational validity to reachability in an abstract reduction system, so the machinery of Section~\ref{sec:ars} applies. If $E$ has an equivalent presentation $\mathcal{R}$ in which every equation is oriented as a rule, and $\mathcal{R}$ is convergent, then $E\models s\approx t$ is decided by comparing normal forms (Corollary~\ref{cor:nf}). \emph{Completion} searches for such an $\mathcal{R}$. Given a reduction order $\succ$ (a well-founded order on terms closed under contexts and substitutions), the Knuth--Bendix procedure maintains a pair $(E_i,\mathcal{R}_i)$ of unoriented equations and oriented rules, starting from $(E,\varnothing)$, using the inferences
\begin{itemize}[nosep]
\item \textsf{Orient}: move $s\approx t$ from $E_i$ to $\mathcal{R}_i$ as $s\to t$ if $s\succ t$ (or as $t\to s$ if $t\succ s$);
\item \textsf{Deduce}: add $s\approx t$ to $E_i$ when $\langle s,t\rangle$ is a critical pair of $\mathcal{R}_i$;
\item \textsf{Simplify}: rewrite a side of an equation, or the right-hand side of a rule, by $\mathcal{R}_i$;
\item \textsf{Delete}: remove an equation $s\approx s$.
\end{itemize}

\begin{theorem}[Correctness of completion]
\label{thm:completion}
\cited{Bachmair, Dershowitz, Hsiang 1986} If a run of the procedure is fair, in the sense that every critical pair between rules that persist from some stage onward is eventually deduced, and the run never meets an equation that is unorientable by $\succ$ and ends with no unoriented equations, then the limit system $\mathcal{R}_\infty$ is convergent and $\leftrightarrow_{\mathcal{R}_\infty}^{*}$ coincides with $\leftrightarrow_E^{*}$. The procedure can fail on an unorientable equation, and can run forever.
\end{theorem}

Completion is a fairness-governed search in the spirit of Chapter~\ref{ch:search} (it deletes as well as adds, so Definition~\ref{def:derivation} does not apply literally), and the discipline of non-certification applies to it verbatim. If completion fails or does not terminate, nothing follows about the decidability of the word problem for $E$, which may be decidable by other means. If it succeeds, the result is a canonizer whose assurance rests on two ingredients: termination of $\mathcal{R}_\infty$ under $\succ$, and the joinability of its critical pairs, both of which a checker can verify afterwards without trusting the run.

\section{Normal forms as a canonical residue}
\label{sec:nf-residue}

Let $\mathcal{R}$ be convergent. The map $\NF:T\to T$ from terms to normal forms is idempotent, sends every term to an equal term, and identifies exactly the $\leftrightarrow^{*}$-equivalent terms (Corollary~\ref{cor:nf}). It is the extreme case of the compression studied in this book: a whole equivalence class, together with every path inside it, is replaced by one representative.

For a task, understood as in Proposition~\ref{prop:factor} of Chapter~\ref{ch:proof-systems} as a function on the full \emph{reduction sequence} rather than on its endpoint, Proposition~\ref{prop:strategy-independence} says that all sequences from $s$ lie in a single fibre of the output map.

\begin{proposition}[What normalization serves and what it cannot]
\label{prop:nf-sufficient}
Let $\mathcal{R}$ be convergent and let $\varrho$ map a maximal reduction sequence from $s$ to its last term.
\begin{enumerate}[label=(\roman*),nosep]
\item The task ``decide whether $s\leftrightarrow^{*}t$'' factors through $\varrho$.
\item In Example~\ref{ex:strategies}, the task ``report the number of rewrite steps used'' does not factor through $\varrho$, and neither does ``report whether the rule $a\to b$ was applied before $h(x)\to d(x,x)$''.
\end{enumerate}
\end{proposition}

\begin{proof}
(i) $s\leftrightarrow^{*}t$ if and only if $\NF(s)=\NF(t)$ by Corollary~\ref{cor:nf}, and $\NF(s)=\varrho(r)$ for every maximal sequence $r$ from $s$. So the answer is a function of the endpoints. (ii) The two sequences from $h(a)$ in Example~\ref{ex:strategies} have the same endpoint $d(b,b)$ and different lengths, and in the second the rule $a\to b$ is applied after $h(x)\to d(x,x)$ while in the first it is applied before. By Proposition~\ref{prop:factor} neither task factors through $\varrho$.
\end{proof}

The proposition is trivial to state and it carries the central message of Part~I. A representation is not sufficient or insufficient in the abstract; it is sufficient \emph{for a declared task}. Normalization is sufficient for equality and for nothing that depends on how equality was reached. Resolution proofs (Chapter~\ref{ch:resolution}) are sufficient for checking and for nothing that depends on what was tried. Part~II shows that SAT solvers, SMT solvers, and kernels exhibit the same pattern at industrial scale, and Part~III turns the pattern into a theory.

\begin{exercise}
Show that the converse of Newman's lemma fails: exhibit a confluent system that is not terminating, and a locally confluent one that is neither terminating nor confluent.
\end{exercise}

\begin{exercise}
Verify by hand the termination argument of Example~\ref{ex:addition} for $\mathit{add}(s(0),s(0))$ and compute its normal form by both innermost and outermost reduction.
\end{exercise}

\begin{exercise}
For the group system of Example~\ref{ex:groups}, compute the normal forms of $(a\cdot b)^{-1}\cdot a$ and of $b^{-1}$, and decide whether they are equal in every group.
\end{exercise}
